% ------------------------------------------------------------------------------
% Chapter 3
% ------------------------------------------------------------------------------
\chapter{Research Methodology} % enter the name of the chapter here
\label{cha:ch3_methodology} % enter the chapter label here (for cross referencing)

% ------------------------------------------------------------------------------
% Introduction
% ------------------------------------------------------------------------------
\section{Introduction}
\label{sec:Introduction}
This chapter identifies and justifies the methodological approaches used to examine the impact of organizational culture and leadership on the return on investment (ROI) of AI coding assistants in software engineering teams. Given the complexity of the constructs to be studied, culture, leadership, intensity of adoption, software engineering performance, and ROI, it is imperative to ensure that the research design adopted is methodologically sound to ensure validity and generalisability. According to \textcite{Creswell2018}, the research design should match the research question and the underlying theory. Since this research intends to test relationships between constructs and assess mediation effects, a quantitative explanatory research design is considered appropriate. 

\section{Research Philosophy}
\label{sec:Research Philosophy}
This research will adopt a post-positivist epistemology. Post-positivists assume that it is possible to examine social phenomena through structured observation. At the same time, they realize that measurement will never be perfect \autocite{Creswell2018}.

\subsection {Post-positivist theory suits this research because it will: }
\label{sec:Post-positivist theory suits this research because it will}
\begin{itemize}
	\item \textbf{Test hypothesized relationships} 
\end{itemize}
Testing hypothesized relationships refers to an evaluation of how well the hypothesized relationships outlined in the conceptual framework are supported by data. In quantitative research, hypotheses are normally developed based on theoretical frameworks available in literature. The hypotheses are then tested using data collected. This allows researchers to evaluate if there is a relationship between various variables, such as organisational culture, leadership style, adoption of AI, and return on investment. Testing of hypotheses is an essential part of quantitative research. This is because it allows researchers to evaluate theoretical frameworks. This process also determines if data is supporting hypothesized relationships or not \autocite{Creswell2018}.

\begin{itemize}
	\item \textbf{Measure latent variables} 
\end{itemize}
Latent constructs refer to concepts or ideas that are abstract and cannot be measured directly, although several measures can be employed for the assessment of the construct. In the study, the constructs such as organizational culture, leadership styles, and the effectiveness of AI implementation can be referred to as latent constructs since they represent complex behaviours in the organization and not directly measurable constructs or variables. Quantitative research measures the constructs through structured survey questions and scales that have been validated for the assessment of the construct in question \autocite{Hair2019}.

\begin{itemize}
	\item \textbf{Estimate causal processes} 
\end{itemize}
To estimate the various causal pathways, it is important to consider the impact of one variable on another through a theoretical framework. In the research, the various causal pathways define the impact of organizational culture and leadership on the effective adoption of AI coding assistants, which eventually influences productivity and return on investment outcomes. To analyse the various relationships between the variables, the structural equation model is often used to examine the various relationships simultaneously, where the direct and indirect effects of the variables on each other can be established. This model provides a better understanding of the various relationships between the variables, as well as the theoretical framework through which the relationships operate \autocite{Hair2019}.

\begin{itemize}
	\item \textbf{Generate generalizable findings} 
\end{itemize}
One of the key goals of quantitative research is to ensure the results obtained are generalizable to a larger population outside the specific sample under consideration. By conducting the research on a sufficiently large sample of software engineering professionals, it is possible to generalize to a larger population of organizations. Statistical techniques help to estimate the likelihood of the observed associations being real rather than coincidental. In this sense, quantitative research provides a framework to produce evidence useful to organizations and the academic sphere \autocite{Bryman2016}.

The use of structured surveys and objective performance measures aligns with the post-positivist school of thought that emphasizes the importance of measurement \autocite{Hair2019}.

\section{Research Design}
\label{sec:Research Design}
\subsection {Quantitative Explanatory Design: }
\label{sec:Quantitative Explanatory Design}
The study will utilize a cross-sectional quantitative design to test hypothesized relationships between organizational culture, leadership practices, successful AI adoption, productivity, and ROI.

Quantitative methods are considered appropriate to examine relationships between variables and to test theoretical models \autocite{Bryman2016}. The study will utilize the methodological strengths of SEM to examine relationships between variables simultaneously, which aligns with the study's objectives \autocite{Hair2019}.

Since the conceptual framework proposes mediation relationships, the study will take advantage of the methodological strengths of SEM to examine relationships between variables simultaneously, which will allow the estimation of the direct, indirect, and total effects between the variables.

\section{Population and Sampling}
\label{sec:Population and Sampling}
\subsection {Target Population: }
\label{sec:Target Population}
The study focuses on software engineering teams within large organizations in South Africa that are using AI coding assistants.

\subsection {Sampling Strategy: }
\label{sec:Sampling Strategy}
The sampling method adopted is purposive sampling, which focuses on organizations that meet the following requirements:
\begin{itemize}
	\item \textbf{The organizations must be using AI coding assistants.} 
	\item \textbf{The organizations must have been using AI coding assistants for at least six months.} 
	\item \textbf{The organizations must have measurable software engineering metrics.} 
\end{itemize}

For the study, the participants will be software developers and team lead within the organizations, using the stratified sampling method.

For the study, the guidelines indicate that the minimum number of participants required to be included in the study using the structural equation modelling method must be more than 200 to ensure the stability of the model \autocite{Hair2019}. Therefore, the study will aim to recruit between 250 and 300 participants.

\section{Measurement and Operationalization}
\label{sec:Measurement and Operationalization}
Latent constructs require measurement scales to be developed to assess construct validity \autocite{DeVellis2022}.

\subsection {Organizational Culture: }
\label{sec:Organizational Culture}
Assessed via adapted and validated scales for:

\begin{itemize}
	\item \textbf{Psychological Safety} 
\end{itemize}
The term \textbf{psychological safety} refers to the level of comfort felt by the members of a given group to experiment, question, and contribute their ideas without the fear of being criticized. In the case of software engineering teams where there is the use of AI coding assistants such as GitHub Copilot, the level of psychological safety allows the developers to question the suggestions given by the tool. When the developers feel safe to experiment with the tool, they will have the ability to become proficient in the use of the tool to effectively utilize the generative AI. According to the literature on the use of generative AI in knowledge work, there is a significant impact on productivity when the workers engage with the tool collaboratively \autocite{Noy2023}.

\begin{itemize}
	\item \textbf{Learning Orientation} 
\end{itemize}
\textbf{Learning orientation} is the term used to describe the organization's emphasis on continuous learning and skill development. As generative AI technologies are integrated into the organization, the need arises for the development of skills by the developers to enable them to efficiently work with the generative AI technologies. Organizations that support a learning orientation are more likely to successfully integrate generative AI technologies into their work processes. Empirical research conducted on the role of generative AI technologies in professional settings has revealed that workers can experience significant improvements in the efficiency of their tasks by working with generative AI technologies \autocite{Brynjolfsson2025}. 

\begin{itemize}
	\item \textbf{Knowledge Sharing} 
\end{itemize}
The term \textbf{knowledge sharing} refers to the sharing of information and expertise by the members of a team. In the context of a team that utilizes AI coding assistants, the sharing of information about prompts, debugging, and the best practices for utilizing the AI code tends to occur. This sharing of information tends to facilitate the ability of the team to utilize generative AI more efficiently. Empirical research concerning the practical utilization of AI coding assistants tends to reveal that the members of a development team rely on the sharing of information to understand the efficient utilization of the AI code \autocite{Sergeyuk2025}.

\begin{itemize}
	\item \textbf{Digital Openness} 
\end{itemize}
The concept of openness to digital change refers to the willingness of organizations and their human resources to accept new digital technologies. Organizations that exhibit openness to digital change are more likely to accept generative AI tools and alter their processes accordingly. In software engineering, openness to digital change is likely to promote the adoption of AI coding assistants by software developers and encourage them to explore new ways of improving their productivity. 

Likert scales can be used to measure perceptions and attitudes \autocite{Bryman2016}.

\subsection {Leadership Practices: }
\label{sec:Leadership Practices}
\textbf{Measured via:}
\begin{itemize}
	\item \textbf{Executive Sponsorship} 
\end{itemize}
\textbf{Executive sponsorship} refers to a situation in which top leaders of an organization publicly support and promote the adoption and implementation of generative AI technologies. In other words, when top leaders publicly support and promote AI coding assistants and other generative AI technologies, it sends a message to other workers that these technologies are of strategic importance and that they are aligned with overall organizational objectives. Research studies that have empirically examined generative AI adoption in knowledge-intensive tasks have shown that when organizations publicly support and promote AI technologies, workers are more likely to engage with AI technologies and explore new ways of task execution \autocite{Noy2023}. In a similar vein, other studies that have empirically examined AI adoption in organizations have shown that executive sponsorship is important in ensuring that AI adoption is effectively coordinated and that workers are able to adopt new technologies with more confidence \autocite{Brynjolfsson2025}.

\begin{itemize}
	\item \textbf{Governance Clarity} 
\end{itemize}
\textbf{Governance clarity} relates to the presence of explicit organisational policies and guidelines on the application of AI technologies. For instance, generative AI tools such as Copilot have salient governance implications on intellectual property rights, security, and software quality assurance. As such, organisations need to have well-defined governance structures to provide a framework on the application of AI-generated code. In addition, governance clarity would help organisations develop accountability mechanisms on the application of AI technology. Studies on the governance of AI have established that organisations with well-defined structures on the application of AI technology are better placed to manage technology risks while reaping the associated productivity benefits \autocite{Hradecky2022}. As such, governance clarity is essential in the sustainable application of generative AI tools in software development environments.

\begin{itemize}
	\item \textbf{Resource Support} 
\end{itemize}
\textbf{Resource support} refers to the provision of financial, technical, and human resources required to support the successful deployment of AI technologies. The deployment of AI coding assistants requires investments in training developers, integrating coding assistants into development environments, infrastructure support, and monitoring. Failure to support these resources may prevent organisations from reaping productivity benefits associated with coding assistants. Empirical research on organisational investments in AI has suggested that organisations that support investments in training developers and digital infrastructure are more likely to achieve performance benefits from investing in AI. \textcite{Babina2024} argue that sufficient resource support is crucial to ensure that AI coding assistants are deployed productively.

\begin{itemize}
	\item \textbf{Metric Alignment} 
\end{itemize}
The issue of \textbf{metric alignment} refers to the alignment of organisational performance metrics with organisational outcomes of AI adoption. In terms of the adoption of AI coding assistants, organisational outcomes are likely to focus on improvements to productivity, reductions to cycle times, and improvements to software quality. Therefore, it is necessary to ensure that these outcomes are aligned with performance metrics such as deployment frequency, defect rates, and development throughput. When performance metrics are aligned with organisational outcomes of AI adoption, employees are more likely to leverage the adoption of AI to achieve organisational outcomes. Research on generative AI productivity has revealed that organisational outcomes become tangible when performance indicators are aligned with AI usage and task efficiencies \autocite{Noy2023}.

\subsection {Effective Adoption: }
\label{sec:Effective Adoption}
\textbf{Effective adoption} refers to the measure of the level to which the coding assistants based on artificial intelligence are assimilated into the software engineering work processes and are used in a productive manner by the software engineers. Adoption of artificial intelligence coding assistants is more than just the installation of the tool, as it is related to the interaction of the software engineers with the tool based on artificial intelligence, the level of trust placed on the tool, and the integration of the tool into the work processes. Empirical studies on artificial intelligence coding assistants have revealed the potential of the tool to increase the productivity of the software engineers, though the level of productivity increase is related to the assimilation of the tool into the work processes \autocite{Brynjolfsson2025}. There are various parameters that can be used to measure the assimilation of artificial intelligence into the work processes of the software engineers.

\textbf{Assessed via:}

\begin{itemize}
	\item \textbf{Usage Intensity} 
\end{itemize}
\textbf{Usage intensity} refers to the frequency of interaction with the coding assistant tool by the developer. Developers who frequently interact with the tool are more likely to become proficient in the use of the tool and to find ways to increase their productivity. It is likely that a high level of usage intensity implies the integration of the tool into the daily work of the developer. Research on the productivity of generative AI in knowledge work has shown that high usage frequency of the tool results in increased productivity and task efficiency \autocite{Noy2023}.

\begin{itemize}
	\item \textbf{Workflow Integration} 
\end{itemize}
The term \textbf{workflow integration} refers to the level of integration of artificial intelligence coding assistants into existing software development processes. This includes integration into the development environment, code review processes, and collaborative programming. It has been established that artificial intelligence coding assistants that integrate well into the existing software development processes have a high probability of being embraced by software developers. Studies on artificial intelligence work systems have established that increased productivity is achieved through the integration of artificial intelligence into existing work processes, as opposed to disrupting work processes \autocite{Brynjolfsson2025}.

\begin{itemize}
	\item \textbf{Trust Calibration} 
\end{itemize}
\textbf{Trust calibration} refers to the ability of developers to trust appropriately the output of an AI while maintaining stringent evaluation. The developers need to understand when to trust appropriately the output of an AI-generated suggestion. Excessive trust in an AI-generated output can lead to errors, while excessive distrust can prevent developers from attaining productivity gains. Studies on collaboration between humans and AI have demonstrated that success depends on building a calibrated trust between humans and an AI, which is a balance between trust and suspicion \autocite{Noy2023}.

\begin{itemize}
	\item \textbf{Feature Utilization Depth} 
\end{itemize}
The utilization depth of features refers to the extent to which developers exploit the advanced capabilities of AI coding assistants, as opposed to using only basic coding features. There are several generative AI tools that provide various features, such as coding, debugging, documentation, and testing. The utilization of a wider range of features is associated with a higher level of productivity. Research studies on AI-assisted software development have shown that a deeper interaction with AI capabilities is associated with improved development and problem-solving efficiency \autocite{Brynjolfsson2025}.

\section{Data Collection}
\label{sec:Data Collection}
The data will be collected through a structured online questionnaire. Online questionnaires are commonly used in organizational and information systems studies due to their effectiveness in collecting data from dispersed samples. The questionnaire will contain closed-ended questions that will be measured using a Likert scale to measure constructs such as organizational culture, leadership practices, successful adoption of AI coding assistants, and organizational outcome.

The questionnaire will be shared through various networks and platforms, such as professional networks, software engineering networks, and technology-oriented platforms. The study will be entirely voluntary, and consent will be sought before data is collected.
The questionnaire will be derived from various measurement scales that have been used in other studies. Using measurement scales that have been used in other studies will ensure that reliable data is collected. The study will measure constructs such as leadership support, organizational culture, AI adoption behaviour, and outcome.

Research studies on AI adoption have shown that AI tool adoption is not entirely dependent on AI technology. The adoption of AI technology also depends on how individuals and organizations adopt and use AI technology. \textcite{Brynjolfsson2025} argue that AI tool adoption is a function of how individuals and organizations adopt and use AI technology. In that case, the study will measure both behavioural and organizational dimensions of AI adoption.

\section{Data Analysis}
\label{sec:Data Analysis}
For this study, data analysis will be conducted using available open-source statistical software. Popular statistical software that is used in scholarly research because they are costless, transparent, and capable of carrying out sophisticated statistical operations such as factor analysis and structural equation modelling. 

The analysis will be conducted in several stages.

The first step in data analysis is carrying out descriptive statistics. Descriptive statistics are used in this study to describe the sample characteristics. In this process, statistics such as means, standard deviations, and frequency distributions are computed. Descriptive statistics are important in data analysis because they help in identifying any missing values and outliers in a dataset before carrying out more sophisticated statistical operations.

The second step in data analysis is conducting tests of \textbf{reliability and validity}. These tests are important in assessing the measurement model used in this study. In this process, Cronbach’s alpha is computed. Cronbach’s alpha is a statistical measure used in assessing the reliability of a set of items in a survey aimed at measuring concepts such as organizational culture, leadership, and effective adoption of AI coding assistants. In organizational and behavioural studies, a Cronbach’s alpha value above 0.70 is considered acceptable \autocite{Hair2019}.

To ensure that the survey items reflect the theoretical constructs embedded within the conceptual framework, confirmatory factor analysis (CFA) will be conducted. This will establish the measurement validity of the research instrument. This is particularly important when researching emergent technologies, such as generative AI. This is because emergent technologies often involve various organisational and behavioural factors when integrating new technologies with work processes \autocite{Mikalef2021}.

Third, in testing the relationships identified in the conceptual framework, this study would use structural equation modelling (SEM). This is because SEM is appropriate in this study in that it would enable this researcher to examine more than one relationship between variables simultaneously. In addition, SEM would allow this researcher to assess both direct and indirect effects in a study involving interdependencies between organizational factors and technology outcomes, which is common in information systems and technology adoption studies (Hair et al., 2019).

For this research, structural equation modelling (SEM) will be used to investigate how organisational culture and leadership practices influence the adoption of AI coding assistants in software engineering teams and whether the adoption of AI coding assistants results in improved productivity and ultimately ROI from AI technology adoption in software engineering teams. Literature on generative AI technology suggests that there is significant scope for improving productivity in knowledge work activities with the adoption and integration of AI technology \autocite{Noy2023}. Similarly, literature on the adoption and integration of AI technology in an organisational context suggests that improvements in productivity due to AI technology adoption and integration are dependent on how effectively individuals and organisations utilise and interact with AI technology \autocite{Brynjolfsson2025}.

In this study, structural equation modelling is used to examine how organisational culture and leadership practices influence effective adoption of AI coding assistants by software engineering teams, as well as whether such effective adoption leads to productivity improvements and ultimately organisational value and return on investment. Previous research on generative AI technologies has pointed out that AI technologies have considerable potential for increasing productivity in tasks such as knowledge work, if users of these technologies are able to effectively integrate these technologies into their work processes \autocite{Noy2023}. Related research in organisational settings has pointed out that productivity enhancements from AI technologies depend not only on the technologies employed but also on users’ interactions with these technologies and organisational support for these technologies’ adoption \autocite{Brynjolfsson2025}.

To test the mediating effects of organizational culture and leadership practices in this conceptual framework, bootstrapping will be used as a statistical method. Bootstrapping is a resampling procedure where a dataset is repeatedly resampled to test the statistical significance of indirect relationships between variables. This method is popular for mediation analysis because it allows for more reliable statistical inferences by providing confidence intervals for indirect effects \autocite{Hair2019}.
This research aims to contribute to an understanding of how organisational culture and leadership practices influence the adoption of AI coding assistants and how this ultimately impacts return on investment from AI technology adoption in software engineering teams.

\section{Ethical Considerations}
\label{sec:Ethical Considerations}
Ethical issues are an integral component of this research. The survey is voluntary, and the respondents will be informed of the purpose of the research before they can respond to the survey questions. They will be made aware that the responses will be anonymous and confidential.

Personal identifiers will not be collected from the respondents. The information will be used only for academic purposes and will be stored in a secure environment. This ensures that the research adheres to the ethical principles of conducting social science research.

\section{Summary}
\label{sec:Summary}
This chapter presents the methodological approach used in the study. A quantitative approach is employed to investigate the relationship between return on investment and organisational culture, leadership practices, and the effective use of AI coding assistants. The data will be collected using a structured survey of software engineering professionals, and statistical analysis methods, such as structural equation modelling, will be used to test the conceptual framework.

Using this methodology, the research aims to produce empirical evidence on the impact of organisational and leadership factors on the value created by generative AI tools in software engineering.